% ---------------------------------------------
% 第1章 電力市場の制度的背景
% ---------------------------------------------

\chapter[電力市場の制度的背景]{電力市場の制度的背景と構造的特徴}

本章では、日本の新電力（小売電気事業者）が直面する制度的・構造的環境を、国際的潮流と国内制度の双方から体系的に整理する。世界的には、脱炭素化や再生可能エネルギーの急速な普及、地政学リスクの高まり、燃料価格の変動など、電力市場の不確実性を増幅させる要因が複合的に作用している。これらの変化は、電力市場の自由化と市場設計の再構築を促し、事業者間のリスク分担構造を大きく変えてきた。

日本においても、2016年の電力小売全面自由化を契機に、発電・送配電・小売の分離、卸電力市場の整備、需給調整市場の高度化など、制度改革が段階的に進展している。しかし、再エネの変動性、燃料価格の高騰、インバランス料金制度、託送料金制度など、日本特有の制度的特徴が新電力の経営リスクを増幅させている点は十分に研究されていない。

本章は、こうした制度的背景を踏まえつつ、新電力のビジネスモデル、卸電力市場の構造、リスク要因の体系を整理し、最終的に「倒産リスク」という本論文の中心テーマへと論理的に接続することを目的とする。

\vspace{1em}

本論文は日本の新電力市場を主たる分析対象とするが、倒産リスクの構造的特徴をより普遍的に理解するため、
を目的とする。  制度構造の異なる英国市場も比較対象として扱う。\footnote{英国市場を比較対象とする理由は二点ある。第一に、英国は1990年代から自由化を先行しており、価格キャップ制度や SLR（Supplier of Last Resort）制度など、小売事業者のキャッシュフローに直接影響する制度が整備されている点で、日本より制度的負荷が大きい。第二に、倒産ステータスが \textit{Active → In administration → Liquidation → Dissolved} と段階的に進行し、倒産日としては \textit{Administration 開始日} を用いるのが一般的であるなど、制度的に明確な倒産判定基準が存在するため、実証分析に適している。}
英国は日本より自由化の歴史が長く、
再エネ比率の上昇、燃料価格の変動、価格キャップ制度、SLR（Supplier of Last Resort）制度など、
小売事業者のキャッシュフローに直接影響を与える制度的特徴を有している。

興味深い点は、日本と英国という制度環境の異なる二つの市場において\footnote{日本では法的整理・私的整理の区分が中心である一方、英国では倒産ステータスが \textit{Active}（通常営業）→\textit{In administration}（管財人管理）→\textit{Liquidation}（清算）→\textit{Dissolved}（法人消滅）と段階的に進む。特に \textit{Administration} は債権者保護と事業継続を目的とする制度であり、倒産日としては Administration 開始日を用いるのが国際的に標準的である。この制度差は、倒産リスクの計測やイベントスタディの設計に直接影響する。}
、
\textbf{「財務指標が良好であっても倒産する企業が存在する」}という現象が共通して観察されることである。
すなわち、倒産は単なる財務指標の悪化ではなく、外部ショックと制度的負担が
企業の構造的脆弱性に掛け算的に作用することで急激に顕在化する。

この共通性は、本論文が後章で導入する「四要素モデル」
（収益性 $\mu$、外部ショック感応度 $\sigma$、支援能力 $S$、倒産閾値 $D$）が、
日本市場だけでなく英国市場にも適用可能であることを示唆する。
以下では、こうした構造的特徴を直感的に理解するために、
日本と英国の両市場で実際に観察された「財務指標では説明できない倒産」の具体例を示す。


\begin{itemize}
    \item \textbf{日本の事例（1）：2021年スポット価格高騰期の新電力 A 社}  
    帝国データバンクおよび東京商工リサーチの倒産報告\parencite{TDB2021PowerCrisis,TSR2021PowerCrisis}によれば、
    A社は直前期まで黒字決算を維持していたものの、2021年1月のJEPX価格急騰により
    調達費用が急増し、数週間で資金繰りが破綻した。財務指標では健全と評価されていた企業が、
    外部ショックにより瞬時に倒産へ至った典型例である。

    \item \textbf{日本の事例（2）：自治体支援の有無で明暗が分かれた B 社と C 社}  
    電気新聞の報道および自治体議会資料\parencite{DenkiShimbun2021Crisis,CityCouncil2021Support}によれば、
    B社は自治体との包括契約に基づき短期的な資金支援を受け倒産を回避した。
    一方、C社は財務指標がB社より良好であったにもかかわらず支援主体が存在せず、
    資金繰りが急速に悪化し撤退に追い込まれた。この対照的な結果は、
    財務指標では捉えられない「支援能力」の重要性を示している。

    \item \textbf{英国の事例（1）：価格キャップ制度改定後に倒産した Bulb Energy}  
    英国規制当局 Ofgem の公式記録\parencite{Ofgem2021BulbSLR}によれば、
    Bulb Energy は顧客数180万を抱える大手であり、直前期の財務指標も急激な悪化は見られなかった。
    しかし、価格キャップ制度により調達費用の上昇を料金へ転嫁できず、
    キャッシュフローが急速に悪化し倒産に至った。

    \item \textbf{英国の事例（2）：同じショックでも生存した Octopus Energy}  
    英国政府および Ofgem の公表資料\parencite{Ofgem2021OctopusTakeover,UKGov2021BulbSupport}によれば、
    同時期に同規模の外部ショックを受けた Octopus Energy は、親会社の資本支援と
    ヘッジ契約の厚さにより倒産を回避した。財務指標だけを比較すれば Bulb と大差はなかったが、
    「支援能力」と「ショック感応度」の違いが生存・倒産を分けた。
\end{itemize}


これらの事例は、財務指標だけでは新電力企業の倒産を説明できないことを示している。\footnote{財務指標が良好であっても倒産が発生する背景には、外部ショックの大きさ（価格変動・燃料価格・制度変更）、ショック感応度（ヘッジ比率・調達構造）、支援能力（親会社・自治体・資本市場へのアクセス）、倒産閾値（制度的負担やキャッシュフロー制約）といった複数要因が掛け算的に作用する構造がある。本論文が後章で導入する線形倒産距離 $Z_{\text{lin}}$ は、この構造的要因を Merton 型モデルおよび Bharath \& Shumway（2008）の枠組みから抽象化したものである。}

同じ外部ショックを受けても倒産する企業と生存する企業が存在する背景には、
収益性、外部ショック感応度、支援能力、倒産閾値という複数の構造的要因が
掛け算的に作用するという特徴がある。

本論文は、この「財務指標では説明できない倒産構造」を理論的に形式化するために、
Merton 型モデルおよび Bharath \& Shumway（2008）の倒産距離の構造を抽象化し、
新電力企業に適合する線形倒産距離 $Z_{\text{lin}}$ を構築する。


新電力市場は単なる政策的副産物ではなく、エネルギー転換期における制度設計の成否を映し出す重要な研究対象であることを示す。

本章で整理する制度的背景は、本論文の中心テーマである「新電力企業の倒産リスク」を、企業金融論の普遍的モデル（Merton モデル・Distance to Default）によって分析するための基礎となる。倒産リスクの発現メカニズムは、国や制度を超えて企業の財務構造に内在する普遍的性質に依拠する一方、電力市場特有の制度設計や外生ショックが企業の経営環境に大きな影響を与える。本章では、日本市場を中心に制度的特徴を整理しつつ、後章で扱う英国市場との共通性・相違性を踏まえ、企業の財務悪化がどのような制度環境のもとで進行するのかを理解するための枠組みを提示する。


% ---------------------------------------------
% 1.1 世界のエネルギー情勢と電力市場の変容
% ---------------------------------------------
\section{世界のエネルギー情勢と電力市場の変容}

本節では、世界的なエネルギー転換の潮流と、それに伴う電力市場の構造変化を概観する。脱炭素化や再生可能エネルギーの急速な普及は、電力システムの変動性を高め、従来の集中型電力供給モデルでは対応しきれない新たな需給調整課題を生み出している。また、地政学リスクや国際燃料価格の変動は、自由化された電力市場における価格形成を不安定化させ、小売電気事業者の財務リスクを増幅させている。さらに、需給逼迫の頻発は市場価格の急騰やインバランス料金の高騰を引き起こし、事業者のリスク管理能力の重要性を高めている。これらの国際的潮流は、日本の電力システム改革にも直接的な影響を与えており、2016年以降の新電力を取り巻く制度環境を理解するうえで不可欠である。本節では、こうした世界的な構造変化を体系的に整理し、次節以降で扱う日本の制度改革や新電力のリスク構造の分析に向けた基礎的枠組みを提示する。

なお、欧州の中でも英国は、自由化の歴史が長く、再エネ導入、価格変動、制度コストの構造が日本市場と類似している。英国市場では、価格キャップ制度や SLR（Supplier of Last Resort）制度の変更が小売事業者のキャッシュフローに直接波及し、外部ショック型の倒産が多数発生している。本研究では、日本市場で観察された倒産メカニズムが、制度構造の異なる英国市場でも再現されるのかを検証することで、倒産リスクの構造原理が国境を越えて普遍的に成立するかを明らかにする。


% ---------------------------------------------
% 1.1.1 供給側の構造変化：再エネ拡大と電力システムの変動性
% ---------------------------------------------
\subsection{供給側の構造変化：再エネ拡大と電力システムの変動性}

21世紀以降、世界のエネルギー市場は、化石燃料依存型の供給構造から、脱炭素化・分散化・デジタル化を基軸とする新たなパラダイムへと移行している。特に2015年のパリ協定以降、温室効果ガス排出削減は国際的義務として位置づけられ、各国は再生可能エネルギー（以下、再エネ）の導入拡大を国家戦略の中心に据えるようになった\parencite{IEA2021NetZero,UNFCCC2015Paris}。

再エネの急速な普及は統計的にも顕著であり、IEA によれば、2023 年の世界の再エネ発電容量は 510 GW 増加し、過去最大を記録した。特に太陽光は年間 380 GW 以上増加し、世界の新規電源の 75\% を占める\parencite{IEA2023Renewables}。こうした変動性電源の比率上昇に伴い、欧州では風力発電の出力が 1 日で 200 GWh 以上変動する事例が確認されるなど、需給調整の難易度が急速に高まっている\parencite{IEA2023Electricity}。

一方、老朽火力の休廃止が進む OECD 諸国では、2030 年までに 80 GW の調整力が不足するとの推計もあり、供給側の柔軟性不足は構造的な課題となっている\parencite{OECD2022Security}。このように、再エネ拡大は脱炭素化の推進力であると同時に、電力システムの変動性と不確実性を増幅させる要因となっている。

% ---------------------------------------------
% 1.1.2 政治・地政学リスク：燃料価格と国際情勢の不確実性
% ---------------------------------------------
\subsection{政治・地政学リスク：燃料価格と国際情勢の不確実性}

世界の電力市場を不安定化させているもう一つの要因が、政治・地政学リスクの高まりである。中東情勢の緊張、ロシア・ウクライナ戦争、米中対立などは、化石燃料の供給網に直接的な影響を与え、LNG・石炭・石油の価格を大きく変動させてきた。

2021〜2022 年の国際燃料価格高騰では、欧州の LNG 価格（TTF 指標）が 20 €/MWh から 200 €/MWh 超へと 10 倍以上に急騰し、石炭価格は 3 倍、原油価格も 2 倍に上昇した\parencite{IEA2022GasMarket}。ロシア・ウクライナ戦争の影響により、欧州のロシア産ガス依存度は 40\% から 10\% 以下へ急減し、その代替調達コストが電力価格の上昇圧力となった。

その結果、欧州主要国の電力価格は 2022 年に前年比 2〜4 倍に上昇し、IEA は 2021〜2023 年の間に世界各国で延べ 70 回以上の需給逼迫警報が発令されたと報告している\parencite{IEA2024Electricity}。こうした政治・地政学リスクは、燃料調達力の弱い小売電気事業者にとって重大な外生ショックとなる。

% ---------------------------------------------
% 1.1.3 需要側の構造変化：AI・データセンターによる電力需要の急増
% ---------------------------------------------
\subsection{需要側の構造変化：AI・データセンターによる電力需要の急増}

近年では、電力供給の不確実性に加えて、電力需要側の急拡大が新たな構造的課題となっている。特に AI の急速な普及に伴い、データセンターが世界の電力需要を押し上げる主要因となっている。IEA の最新報告によれば、世界のデータセンターの電力消費は 2024 年の 460 TWh から、2030 年には 1,000 TWh を超え、2035 年には 1,300 TWh に達する見通しであり、わずか 10 年で約 3 倍に増加する\parencite{IEA2025EnergySupplyAI}。さらに、AI 最適化型データセンターに限れば、2030 年までに電力需要が 4 倍以上に増加するとの推計もある\parencite{IEA2025AISurge}。

この規模は、2030 年のデータセンター電力消費が日本全体の年間電力消費量に匹敵する水準であることを意味し、電力システムに対する負荷は極めて大きい。また、米国では、2030 年までの電力需要増加の約半分がデータセンターによるものと見込まれており、AI の普及が電力需要構造を根本的に変えつつある\parencite{IEA2025AISurge}。

このように、AI・デジタル化の進展は、電力供給の変動性だけでなく、需要側の急増という新たな不確実性をもたらしており、各国の電力市場設計や供給力確保策に対して従来以上の柔軟性と投資を要求している。

% ---------------------------------------------
% 1.1.4 電力市場自由化の国際的潮流
% ---------------------------------------------
\subsection{電力市場自由化の国際的潮流}

世界の電力市場自由化は、各国の制度設計やエネルギー政策の違いによって多様な形態をとってきたが、その根底には、前節までに整理した供給側の変動性、地政学リスク、需要側の急増という三つの外生ショックに対応するための制度的枠組みを構築するという共通の目的が存在する。本節では、世界の電力市場自由化を「欧州モデル」「北米モデル」「アジアモデル」「新興国モデル」の四つに類型化し、それぞれの政策的特徴と市場構造を整理する。

\subsubsection{（1）欧州モデル：完全自由化と市場統合}

欧州連合（EU）は、世界で最も進んだ電力市場自由化を実現している地域である。発電・送電・小売の完全分離（アンバンドリング）が徹底され、卸電力市場はスポット・先渡・時間前市場が高度に整備されている。さらに、Nord Pool や EPEX Spot に代表されるように、国境を越えた市場統合が進展し、域内の電力取引が活発化している。

EU の再エネ比率は 2023 年に 44\% に達し、変動性電源の増加に伴い市場価格の変動幅も拡大している。Nord Pool では、日間の価格変動が ±200 €/MWh に達する事例も報告されており、需給調整力の確保が重要な政策課題となっている。また、欧州の小売スイッチング率は 年間 10〜20\% と高く、競争が活発である一方、価格変動リスクが小売事業者に直接転嫁される構造が形成されている。

\subsubsection{（2）北米モデル：部分自由化と州別市場設計}

北米では、州ごとに自由化の度合いが大きく異なる。PJM、CAISO、ERCOT などの ISO/RTO が市場運営を担い、需給調整市場や容量市場が高度に発達している。特にテキサス州（ERCOT）は完全自由化市場として知られ、2021 年の寒波時にはスポット価格が 9,000 \$/MWh に達するなど、極端な価格スパイクが発生した。

また、北米では AI・データセンター需要が電力需要増加の主要因となっており、IEA によれば 2030 年までの電力需要増加の約 50\% がデータセンター由来と推計されている。このように、北米モデルは市場効率性が高い一方、地域差が極端であり、制度設計の違いが小売事業者のリスク構造に大きな影響を与えている。

\subsubsection{（3）アジアモデル：段階的自由化と国家主導}

アジアでは、日本・韓国・台湾などが段階的な自由化を進めている。発電・小売の自由化は進展しているものの、送配電部門は規制下にあり、国家主導の市場設計が特徴である。再エネ導入は政策主導で急速に進んでいるが、市場設計は欧米に比べて未成熟であり、制度変更リスクが大きい。

日本では、2016 年の小売全面自由化以降、新電力のシェアは 約 20\% に達したが、2021 年のスポット価格高騰時には多数の新電力が経営危機に直面した。韓国の再エネ比率は 9\%（2023 年） と低く、台湾も 12\% にとどまるなど、アジア地域は再エネ比率の低さと制度依存の高さが特徴である。

\subsubsection{（4）新興国モデル：国家主導型市場化}

中国・インド・ブラジルなどの新興国では、国家主導で市場化が進められている。中国は世界最大の再エネ導入国であり、2023 年には 300 GW 以上の再エネ設備を追加した。インドは再エネ比率が 30\% に達し、ブラジルは水力依存度が 60\% を超えるなど、各国のエネルギーミックスは大きく異なる。

これらの国では市場規模が大きい一方、価格規制や国家主導の投資が中心であり、自由化の度合いは限定的である。そのため、小売事業者の市場リスクは比較的小さいが、政策依存度が高く、制度変更が事業環境に直接影響を与える。

\subsubsection{（5）四つのモデルの比較と新電力リスクへの含意}

以上の四つのモデルを比較すると、欧州は価格変動リスクが大きく、北米は地域差が極端であり、アジアは制度リスクが大きく、新興国は政策依存度が高いという特徴が明らかとなる。日本はアジアモデルに属し、段階的自由化と国家主導の制度設計のもとで市場が形成されているが、再エネ比率の低さや市場設計の未成熟さから、外生ショックに対して脆弱な構造を持つ。

このように、世界の電力市場自由化は単一のモデルではなく、各国の制度設計とエネルギー政策の違いによって多様な形態をとっており、これらの違いは新電力のリスク構造に直接的な影響を与える。本論文では、こうした国際的潮流を踏まえつつ、日本の新電力が直面する制度的・市場的リスクを分析する。



% ---------------------------------------------
% 1.2 日本の電力システム改革と市場構造
% ---------------------------------------------
\section{日本の電力システム改革と市場構造}

日本の電力システム改革は、1995年の電気事業法改正に端を発し、段階的に自由化が進められてきた\parencite{METI2015Reform}。

\subsection{1995 年以前の日本の電力システム：地域独占と総括原価方式}

1995 年以前の日本の電力システムは、戦後一貫して「地域独占・総括原価方式」を基盤として構築されてきた。1951 年の電気事業再編により、北海道から沖縄までの 10 電力会社が地域独占の形で発電・送電・配電・小売を一体的に担う体制が確立した。この体制は、戦後復興期から高度経済成長期にかけて、電力の安定供給と大規模投資の実行に一定の成果を上げた。

しかし、1980 年代後半以降、以下の構造的課題が顕在化した。

\begin{itemize}
  \item \textbf{競争の欠如}：地域独占により、電力会社間の競争が存在せず、効率化インセンティブが弱かった。
  \item \textbf{料金の硬直性}：総括原価方式により、燃料費・設備投資・人件費などのコストがそのまま料金に転嫁され、国際比較で高い電気料金が問題視された。
  \item \textbf{電源構成の硬直化}：原子力・大規模火力を中心とした集中型電源が主流であり、再エネや分散型電源の導入が進まなかった。
  \item \textbf{需給調整の非効率性}：広域連系線の容量が小さく、地域間融通が限定的で、需給調整の柔軟性が低かった。
\end{itemize}

これらの課題は、1990 年代に入り、国際的な電力自由化の潮流や国内の電気料金高騰を背景に、制度改革の必要性を高める要因となった。

\subsection{世界的潮流が日本に与えた影響：自由化・脱炭素・市場化の圧力}

1990 年代以降、世界の電力市場は、欧米を中心に自由化・市場化が急速に進展した。英国では 1990 年の電力法改正により発電・送電・小売が分離され、競争市場が形成された。米国でも 1990 年代半ばから州単位での電力自由化が進み、ISO/RTO による市場運営が一般化した。さらに、欧州連合（EU）は 1996 年の電力指令を皮切りに、域内市場の統合と完全自由化を段階的に実施した\parencite{IEA2001Reform,EU1996Directive}。

こうした国際的潮流は、日本の電力制度にも大きな影響を与えた。第一に、\textbf{国際比較で高水準の電気料金}が国内産業の競争力を低下させているとの指摘が強まった。1990 年代初頭、日本の産業用電気料金は OECD 諸国の中でも高く、製造業の国際競争力に対する懸念が政策議論の中心となった\parencite{OECD1994EnergyPrice}。

第二に、\textbf{再エネ導入の遅れ}が国際的な批判を招いた。欧州では風力・太陽光の導入が急速に進む一方、日本では制度的制約や電力会社の垂直統合構造により、分散型電源の普及が進まなかった。1992 年の地球サミット以降、温室効果ガス削減が国際的課題となる中、日本の電源構成の硬直性は政策的な弱点として認識された。

第三に、\textbf{市場の透明性と競争の欠如}が国際機関から指摘された。世界銀行や OECD は、日本の電力市場が地域独占・総括原価方式に依存している点を問題視し、競争導入による効率化を提言した\parencite{WorldBank1993PowerSector}。

これらの外部環境の変化に加え、国内でもバブル崩壊後の経済停滞を背景に、産業界から電気料金引き下げと制度改革を求める声が高まった。特に経団連は、1994 年以降、電力自由化を明確に政策提言として掲げ、政府に対して制度改革を強く求めた\parencite{Keidanren1994Proposal}。

このように、1990 年代の日本は、\textbf{国際的潮流（自由化・脱炭素・市場化）と国内の経済構造変化}の双方から、電力制度改革を迫られる状況にあった。これが、1995 年の電気事業法改正へとつながる政策的背景である。

\subsection{電力システム改革の政策形成：経産省・電力会社・有識者会議の役割}

1990 年代の日本における電力システム改革は、単なる制度改正ではなく、政府、電力会社、産業界、有識者会議といった複数のアクターが関与する政策過程の中で形成された。特に、1995 年の電気事業法改正は、国際的潮流と国内の構造的課題を背景に、経済産業省（当時は通商産業省）が中心となって制度設計を主導した点に特徴がある。

まず、\textbf{通商産業省（現・経済産業省）}は、電力制度改革の政策形成において中心的役割を果たした。1990 年代初頭、産業界からの電気料金引き下げ要求が強まる中、通産省は「電気事業制度改革研究会」や「電気事業審議会」を設置し、制度改革の方向性を検討した。これらの会議では、欧米の自由化動向や国内の料金水準を踏まえ、部分的な競争導入が必要であるとの認識が共有された\parencite{METI1994Council}。

次に、\textbf{電力会社（一般電気事業者）}は、地域独占体制の維持を主張しつつも、国際的潮流や国内の批判を受け、一定の制度改革を受け入れざるを得ない状況にあった。特に、総括原価方式の見直しや発電部門への競争導入については慎重姿勢を示したが、通産省との調整を通じて、段階的自由化を前提とした制度改革案が形成された\parencite{Denjiren1995Position}。

また、\textbf{産業界（経団連など）}は、電気料金の国際比較に基づき、電力自由化を強く要求した。経団連は 1994 年以降、電力市場への競争導入を政策提言として掲げ、政府に対して制度改革を促した。これにより、電力制度改革は「産業競争力強化」の文脈でも重要な政策課題として位置づけられた\parencite{Keidanren1994Proposal}。

さらに、\textbf{有識者会議}は、制度改革の方向性を理論的に支える役割を果たした。特に、電力市場の透明性向上、競争導入による効率化、広域的な需給調整の必要性などが議論され、これらの提言が 1995 年の電気事業法改正に反映された\parencite{Expert1995Reform}。

このように、1995 年の電気事業法改正は、通産省が主導しつつも、電力会社、産業界、有識者会議といった複数のアクターが相互に影響を及ぼしながら形成された政策である。改革の基本方針は、\textbf{「競争導入による効率化」「料金の透明性向上」「広域的な需給調整の強化」}であり、これがその後の三段階自由化（発電自由化、小売自由化、送配電分離）へとつながる制度的基盤となった。

\subsection{改正電気事業法と三段階自由化：発電・小売・送配電の再設計}

1995 年の電気事業法改正は、日本の電力システム改革の出発点であり、その後 20 年以上にわたって段階的に進められる自由化の基盤を形成した。本節では、1995 年以降の制度改革を「三段階自由化」として整理し、発電・小売・送配電の各部門がどのように再設計されたかを概観する。

\subsubsection{（1）第一段階：1995 年改正電気事業法による部分自由化}

1995 年の改正では、特定規模需要（いわゆる大口需要家）に対する部分的な小売自由化が導入された。これにより、一定規模以上の工場やオフィスは、一般電気事業者以外の事業者から電力を調達することが可能となった。また、発電事業は「届出制」とされ、新規参入が容易になった。これにより、独立系発電事業者（IPP）が市場に参入し、発電部門に限定的ながら競争が導入された\parencite{METI1995Reform}。

しかし、この段階では送配電部門は依然として地域独占のままであり、競争の範囲は限定的であった。とはいえ、1995 年改正は、垂直統合型の電力システムに初めて競争原理を導入した点で歴史的意義を持つ。

\subsubsection{（2）第二段階：2000 年代の発電自由化と小売自由化の拡大}

2000 年の改正では、発電部門の自由化がさらに進展し、卸電力取引所（JEPX）が設立された。これにより、発電事業者間の競争が市場メカニズムを通じて行われるようになり、価格の透明性が向上した。また、2004 年・2005 年の改正により、小売自由化の対象が段階的に拡大され、特別高圧・高圧需要家の大部分が自由化市場に移行した\parencite{METI2005Reform}。

この時期には、広域的な需給調整の必要性が高まり、広域系統運用の議論が進んだ。再エネ導入の拡大や電力需要の変動性の増大に対応するため、地域間連系線の強化や広域的な系統運用の枠組みが検討されるようになった。

\subsubsection{（3）第三段階：2016 年の小売全面自由化と 2020 年の送配電分離}

2016 年の電気事業法改正により、家庭用を含むすべての需要家が自由化市場で電力会社を選択できるようになった。これにより、小売市場は完全自由化され、多数の新電力が参入した。小売全面自由化は、電力市場における競争を本格化させ、料金メニューの多様化やサービス競争を促進した\parencite{METI2016FullRetail}。

さらに、2020 年には送配電部門の法的分離（アンバンドリング）が実施され、発電・小売と送配電の機能が明確に分離された。これにより、送配電網の中立性が制度的に担保され、競争市場の公平性が強化された。送配電部門は「一般送配電事業者」として位置づけられ、ネットワークの安定運用と中立的な接続サービスの提供が求められるようになった\parencite{METI2020Unbundling}。

\subsubsection{（4）三段階自由化の制度的意義}

以上の三段階自由化により、日本の電力市場は、発電・小売の競争市場と、送配電の規制部門が併存する「部分的競争市場モデル」へと移行した。これは欧州モデルに近い構造であり、競争導入による効率化と、ネットワーク部門の中立性確保を両立させる制度設計である。

しかし、自由化の進展に伴い、価格変動リスクの顕在化、需給逼迫時の市場機能の脆弱性、新電力の経営基盤の弱さなど、新たな課題も浮上した。これらの課題は、次節で検討する新電力のビジネスモデルと密接に関連している。

\subsubsection{自由化後の市場構造の変容：競争導入・価格形成・市場の成熟度}

三段階自由化を経て、日本の電力市場は、発電・小売の競争市場と送配電の規制部門が併存する構造へと移行した。本節では、自由化後の市場構造がどのように変容したかを、競争導入、価格形成、市場成熟度の三つの観点から整理する。

\subsubsection{（1）競争導入：新電力の参入と市場シェアの変化}

2016 年の小売全面自由化により、多数の新電力が市場に参入した。小売電気事業者の登録数は、自由化前の十数社から、2020 年には 700 社を超える規模に拡大した\parencite{METI2020RetailData}。これにより、家庭用市場を含む小売市場に競争が導入され、料金メニューの多様化やサービス競争が進展した。

市場シェアの面では、新電力の販売電力量シェアは、自由化直後の 5\% 程度から、2023 年には約 20\% に達した。特に都市部ではスイッチング率が高く、東京電力管内では家庭用の約 25\% が新電力へ移行したとされる。一方、地方部では依然として旧一般電気事業者のシェアが高く、地域間で競争の進展度合いに差が見られる。

\subsubsection{（2）価格形成：JEPX を中心とした市場価格の変動性}

自由化後の価格形成において重要な役割を果たすのが、卸電力市場（JEPX）である。発電事業者と小売事業者は、スポット市場や先渡市場を通じて電力を取引し、その価格が小売料金や調達コストに直接影響を与える。

JEPX のスポット価格は、再エネ比率の上昇や燃料価格の変動、需給逼迫時の市場環境に応じて大きく変動する。特に 2021 年 1 月の価格高騰では、スポット価格が一時 200 円/kWh を超える異常値を記録し、多くの新電力が調達コストの急騰に直面した\parencite{JEPX2021PriceSpike}。この事例は、自由化市場における価格変動リスクの大きさを象徴している。

また、容量市場や需給調整市場の導入により、電源投資や調整力確保のための新たな価格シグナルが形成されつつあるが、これらの市場はまだ発展途上であり、価格メカニズムの透明性や予見可能性には課題が残る。

\subsubsection{（3）市場成熟度：競争の実効性と制度的課題}

自由化後の市場は一定の競争を実現したものの、その成熟度には限界がある。第一に、旧一般電気事業者の市場支配力が依然として強く、発電・小売の両部門で高いシェアを維持している。特に発電部門では、旧電力会社が依然として 70\% 以上の供給力を保有しており、新電力の競争力を制約する要因となっている\parencite{OCCTO2022Supply}。

第二に、送配電部門の中立性は制度的に確保されたものの、接続ルールや託送料金制度の複雑さが新規参入者にとって障壁となっている。特に再エネ電源の接続制約や出力抑制の問題は、自由化市場の効率性を損なう要因として指摘されている。

第三に、需要家側のスイッチング行動には地域差が大きく、価格感応度の低さや情報不足が競争の実効性を弱めている。家庭用市場では、自由化後も約 70\% が旧一般電気事業者の料金メニューを継続しており、競争の浸透度には限界がある。

以上のように、自由化後の日本の電力市場は、競争導入と市場メカニズムの構築に一定の成果を上げたものの、価格変動リスクの顕在化、旧電力会社の市場支配力、制度的な参入障壁など、成熟した競争市場としての課題を抱えている。これらの課題は、新電力のビジネスモデルや経営リスクに直接的な影響を与えるため、次節で詳細に検討する。


% ---------------------------------------------
% 1.3 新電力のビジネスモデル
% 調達・需給管理・販売・回収の4機能を体系的に整理し、図1で構造を可視化する。
% ---------------------------------------------

\section{新電力のビジネスモデル}

本節では、新電力（小売電気事業者）のビジネスモデルを、調達・需給管理・販売・回収という4つの主要機能に分解し、その構造的特徴を整理する。

% ---------------------------------------------
% 図1（ランドスケープ＋TikZ）
% ---------------------------------------------
\begin{landscape}
\begin{figure}[H]
\centering
\begin{tikzpicture}[
    box/.style={
        rectangle,
        draw,
        rounded corners,
        align=center,
        minimum width=48mm,
        minimum height=12mm,
        font=\small
    },
    smallbox/.style={
        rectangle,
        draw,
        rounded corners,
        align=center,
        minimum width=40mm,
        minimum height=12mm,
        font=\small
    },
    arrow/.style={->, thick}
]

% --- 上段（中央に揃える） ---
\node[box] (policy) at (6,7) {制度設計・規制\\（経産省・OCCTO）};
\node[box] (cost)   at (6,5) {制度コスト\\（需給調整金・容量市場）};

% --- 右上（外部ショック） ---
\node[box] (shock)  at (13,6.5) {外部ショック\\（燃料価格・気象・需給逼迫）};

% --- 左側 ---
\node[box] (gen1)   at (-2,3) {大規模発電事業者\\（例：東京電力・JERA等）};
\node[box] (jepx)   at (-2,1) {卸電力市場（JEPX）\\（スポット・先渡・長期契約）};
\node[box] (fit)    at (-2,-1) {FIT電源\\（固定価格買取制度）};
\node[box] (selfgen)at (-2,-3) {自社電源\\（太陽光・風力・バイオマス等）};

% --- 中央 ---
\node[box] (retailer) at (6,1) {新電力（小売電気事業者）\\（需給管理・インバランス管理）};

% --- 右側（横幅を縮めた顧客ボックス） ---
\node[smallbox] (customer) at (14,1)
{顧客\\（料金プラン・地域サービス・\\再エネメニュー）};

% --- 下段 ---
\node[box] (imbalance) at (6,-2.5) {インバランス料金\\（需給誤差ペナルティ）};
\node[box] (risk)      at (6,-5) {新電力の損益・倒産リスク};

% ============================================================
%  矢印（あなたの指定どおりに完全再現）
% ============================================================

% --- 左側：大規模発電 → JEPX（下向き） ---
\draw[arrow] (gen1.south) -- (jepx.north) node[midway, right] {供給};

% --- 左側：FIT → JEPX（上向き） ---
\draw[arrow] (fit.north) -- (jepx.south) node[midway, right] {市場へ放出};

% --- 左側：自社電源 → 新電力（斜め） ---
\draw[arrow] (selfgen.east) -- (retailer.south west);

% --- 左側：JEPX → 新電力（右向き） ---
\draw[arrow] (jepx.east) -- (retailer.west) node[midway, above] {外部調達};

% --- 上段：制度設計 → 制度コスト（下向き） ---
\draw[arrow] (policy.south) -- (cost.north);

% --- 上段：制度コスト → 新電力（下向き、距離短く） ---
\draw[arrow] (cost.south) -- (retailer.north) node[midway, right] {負担};

% --- 外部ショック → 新電力（斜め） ---
\draw[arrow] (shock.south west) -- (retailer.north east);

% --- 中央：新電力 → インバランス料金（下向き） ---
\draw[arrow] (retailer.south) -- (imbalance.north) node[midway, right] {需給誤差};

% --- 下段：インバランス料金 → 損益・倒産リスク（下向き） ---
\draw[arrow] (imbalance.south) -- (risk.north);

% --- 新電力 ↔ 顧客（双方向） ---
\draw[arrow] (retailer.east) -- (customer.west) node[midway, above] {販売};
\draw[arrow] (customer.west) -- (retailer.east) node[midway, below] {料金回収};

\end{tikzpicture}

\caption{新電力のビジネスモデルの全体構造図}
\label{fig:business_model}
\end{figure}
\end{landscape}

% ---------------------------------------------
% 図1の後の本文（ここが欠落していた部分）
% ---------------------------------------------

% ---------------------------------------------
% 1.3 新電力ビジネスの全体構造：調達・需給管理・販売の三層モデル
% ---------------------------------------------

\subsection{新電力ビジネスの全体構造：調達・需給管理・販売の三層モデル}

新電力（小売電気事業者）のビジネスモデルは、図 1 に示すように、
\textbf{（1）電力の調達}、\textbf{（2）需給管理}、\textbf{（3）販売（小売）}の三つの主要プロセスから構成される。
これらは左から右へと流れる一連のバリューチェーンを形成しており、
新電力はこの三層を統合的に運営することで事業を成立させている。

第一に、\textbf{調達}では、発電事業者からの相対契約、卸電力市場（JEPX）でのスポット・先渡取引、
再生可能エネルギーの固定価格買取制度（FIT）電源など、多様な調達手段を組み合わせる。
調達コストは新電力の収益構造を大きく左右するため、価格変動リスクの管理が重要となる。

第二に、\textbf{需給管理}では、需要家の電力使用量と調達電力量を一致させる必要がある。
需給が一致しない場合、その差分は「インバランス」として精算され、
不足時には高額のインバランス料金が発生する。
このため、需要予測、調達計画、リアルタイムの需給調整が新電力の中核的機能となる。

第三に、\textbf{販売（小売）}では、需要家に対して料金メニューや付加価値サービスを提供し、
顧客獲得と収益確保を図る。
料金メニューには、従量料金、時間帯別料金、セット割、再エネメニューなどがあり、
需要家の多様なニーズに応じたサービス設計が求められる。

この三層モデルには、発電事業者、一般送配電事業者、広域機関（OCCTO）、
卸電力市場（JEPX）、需要家など、多様な主体が関与しており、
それぞれが異なる利害関係を持つ。
新電力は、これらの主体との関係を調整しつつ、
市場価格、制度設計、需給状況といった外部環境の変動に対応しながら事業を運営する必要がある。

本節では、この三層モデルに沿って、新電力のビジネス構造を詳細に検討する。


\subsection{調達：発電事業者・卸電力市場・FIT 電源の位置づけ}

新電力のビジネスモデルにおいて、調達は最も重要な基盤であり、事業収益の大部分を決定する要素である。
図 1 に示すように、調達は主として
\textbf{（1）発電事業者との相対契約}、
\textbf{（2）卸電力市場（JEPX）での取引}、
\textbf{（3）FIT 電源の受け入れ}
の三つの手段から構成される。

\subsubsection{電気という商品の同質性と託送の必然性}

電気は、発電主体や発電方法が異なっても、需要家に届けられる時点では
\textbf{完全に同質的（コモディティ）}な商品である。
すなわち、電気そのものには差別化要素がなく、自由市場における一般的な
商品競争とは異なる構造を持つ。

しかし、電気を需要家に届けるためには、一般送配電事業者が保有する
送電・配電ネットワークを必ず利用する必要がある。
このネットワーク利用行為が「\textbf{託送}」であり、
新電力は託送に対して「\textbf{託送料金}」を支払う。
託送料金は、送配電網の維持・運用に必要なコストを反映したものであり、
小売電気料金の重要な構成要素となる。

したがって、新電力の調達コストは、電源調達費用（相対契約・市場調達・FIT）
に加えて、託送料金を含むネットワーク利用コストを総合的に考慮する必要がある。
この点で、電気の同質性とネットワーク依存性は、新電力ビジネスの
構造的特徴を規定する重要な要素である。

\subsubsection{電気の同質性がもたらす付加価値競争の構造}

電気が完全に同質的な商品であるという特性は、新電力の競争戦略に直接的な影響を与える。
すなわち、需要家に届けられる電気そのものには差別化要素が存在しないため、
新電力は「電気以外の価値」で競争せざるを得ない。

この構造は、一般的な消費財市場における自由競争とは本質的に異なる。
新電力が提示する差別化要素は、電気そのものではなく、以下のような
\textbf{付加価値サービス}に依存している。

\begin{itemize}
  \item \textbf{環境価値の提供}：実質再エネメニュー、CO\textsubscript{2} 排出量の可視化
  \item \textbf{料金以外の経済的メリット}：ガス・通信とのセット割、ポイント還元
  \item \textbf{IT を活用したエネルギーマネジメント}：スマートホーム連携、需要最適化サービス
\end{itemize}

このように、電気の同質性とネットワーク依存性という構造的特徴は、
新電力が「電気以外の価値」で競争する市場環境を生み出している。
これは、後述する販売戦略（1.3.4）における付加価値サービスの重要性とも密接に関連している。

\subsubsection{（1）発電事業者との相対契約}

相対契約とは、新電力が発電事業者（IPP、再エネ事業者、旧一般電気事業者など）と
個別に締結する電力供給契約である。相対契約は、価格の安定性と供給量の確実性が高い点で、
新電力にとって重要な調達手段である。

相対契約の特徴は以下の通りである。

\begin{itemize}
  \item \textbf{価格の固定性}：一定期間の固定価格契約が多く、JEPX の価格変動リスクを回避できる。
  \item \textbf{供給量の確実性}：契約量が事前に確定するため、需給管理の基礎となる。
  \item \textbf{電源属性の選択}：再エネ電源、火力電源など、電源の種類を選択できる。
\end{itemize}

一方で、相対契約は長期契約が多く、柔軟性に欠けるという課題もある。
市場価格が下落した場合でも契約価格が維持されるため、調達コストが割高になる可能性がある。

\subsubsection{（2）卸電力市場（JEPX）での取引}

卸電力市場（Japan Electric Power Exchange, JEPX）は、
新電力が電力を調達する主要な市場であり、スポット市場、先渡市場、時間前市場など複数の市場で構成される。
JEPX の価格は需給状況、燃料価格、再エネ発電量などに応じて変動し、
新電力の調達コストに直接影響を与える。

JEPX の特徴は以下の通りである。

\begin{itemize}
  \item \textbf{スポット市場}：翌日分の電力を 30 分単位で取引する市場で、価格変動が最も大きい。
  \item \textbf{先渡市場}：1 か月〜1 年先の電力を固定価格で取引する市場で、価格安定化に寄与する。
  \item \textbf{時間前市場}：当日内の需給調整に利用される市場で、インバランス回避に重要。
\end{itemize}

新電力は、相対契約だけでは不足する電力量を JEPX で補うことが多く、
JEPX の価格変動は新電力の収益に大きな影響を与える。
特に 2021 年 1 月の価格高騰のような事例では、
調達コストが急騰し、多くの新電力が経営危機に直面した。

\subsubsection{（3）FIT 電源の位置づけ}

FIT（Feed-in Tariff）は、再生可能エネルギーを固定価格で買い取る制度であり、
新電力は「FIT 電気」を国が定めるルールに従って受け入れる。
FIT 電源は、発電事業者が一般送配電事業者に売電し、
その電気を新電力が「FIT 電気」として受け取る仕組みである。

FIT 電源の特徴は以下の通りである。

\begin{itemize}
  \item \textbf{価格は固定}：FIT 買取価格は国が決定し、市場価格とは独立している。
  \item \textbf{環境価値は非付帯}：FIT 電気には環境価値が付帯しないため、非化石証書を別途購入する必要がある。
  \item \textbf{調達量は不確実}：太陽光・風力など変動性電源が多く、供給量が天候に左右される。
\end{itemize}

FIT 電源は価格が安定している一方、供給量の不確実性が高く、
需給管理の難易度を高める要因となる。

\subsubsection{（4）調達ポートフォリオとリスク管理}

新電力は、相対契約、JEPX、FIT 電源を組み合わせて調達ポートフォリオを構築する。
その目的は、以下の三つのリスクを最小化することである。

\begin{itemize}
  \item \textbf{価格リスク}：JEPX の価格変動により調達コストが急増するリスク。
  \item \textbf{量的不確実性}：FIT 電源や再エネ電源の出力変動による供給不足リスク。
  \item \textbf{インバランスリスク}：需給不一致により高額のインバランス料金が発生するリスク。
\end{itemize}

調達ポートフォリオの最適化は、新電力の収益性と経営安定性を左右する中核的な戦略であり、
次節で述べる需給管理と密接に関連している。


\subsection{需給管理：インバランス制度と広域機関の役割}

新電力のビジネスモデルにおいて、需給管理は調達と並ぶ中核的機能であり、
事業の安定性と収益性を左右する最も重要なプロセスである。
需給管理とは、需要家が実際に使用する電力量と、新電力が調達した電力量を
リアルタイムで一致させることであり、日本の電力システムでは
「計画値同時同量」の原則として制度化されている。
本節では、需給管理の仕組みと、インバランス制度および広域機関（OCCTO）の役割を整理する。

\subsubsection{（1）計画値同時同量と需要予測の重要性}

新電力は、需要家の電力使用量を事前に予測し、その予測に基づいて調達計画を作成する。
この予測値（計画値）と実際の使用量（実績値）が一致することが求められ、
一致しない場合、その差分は「インバランス」として精算される。

需要予測は、以下の要素を考慮して行われる。

\begin{itemize}
  \item 過去の需要実績（曜日・季節・気温など）
  \item 需要家の属性（家庭用・法人用、業種、稼働パターン）
  \item 天候予測（特に夏季・冬季の気温）
  \item 再エネ電源の出力予測（FIT・FIP 電源を保有する場合）
\end{itemize}

需要予測の精度が低い場合、インバランスが増加し、経営リスクが高まる。

\subsubsection{（2）インバランス制度：需給不一致の精算メカニズム}

インバランスとは、計画値と実績値の差分に対して課される精算であり、
不足時には高額のインバランス料金が発生する。
インバランス料金は、一般送配電事業者が需給調整のために調達した電力のコストを反映しており、
市場価格が高騰している場合には極めて高額となる。

インバランス制度の特徴は以下の通りである。

\begin{itemize}
  \item \textbf{不足インバランス}：調達量が需要を下回った場合に発生し、最も高額。
  \item \textbf{余剰インバランス}：調達量が需要を上回った場合に発生し、マイナス精算となることもある。
  \item \textbf{市場連動性}：インバランス料金は市場価格（JEPX）や需給逼迫状況に応じて変動する。
\end{itemize}

2021 年 1 月の価格高騰時には、インバランス料金が 200 円/kWh を超える事例も発生し、
多くの新電力が巨額の損失を計上した。
この事例は、需給管理の失敗が新電力の経営に致命的な影響を与えることを示している。

\subsubsection{（3）広域機関（OCCTO）の役割}

広域機関（OCCTO：電力広域的運営推進機関）は、
日本の電力システムにおける需給管理の中核的な制度主体であり、
以下の役割を担っている。

\begin{itemize}
  \item \textbf{需給計画の集約}：すべての小売事業者・発電事業者の計画値を集約し、需給バランスを監視する。
  \item \textbf{広域的な需給調整}：地域間連系線を活用し、需給逼迫時に広域的な調整を行う。
  \item \textbf{インバランス精算のルール策定}：インバランス料金の算定方法や精算ルールを定める。
  \item \textbf{供給力評価}：各事業者の供給力を評価し、需給逼迫リスクを事前に把握する。
\end{itemize}

OCCTO は、自由化市場における「需給の番人」として機能しており、
新電力は OCCTO が定めるルールに従って需給管理を行う必要がある。

\subsubsection{（4）一般送配電事業者の役割}

一般送配電事業者は、送配電網の運用と需給調整を担う主体であり、
新電力の需給管理に直接関与する。

主な役割は以下の通りである。

\begin{itemize}
  \item \textbf{需給実績の計測}：スマートメーターを通じて需要家の実績値を計測する。
  \item \textbf{インバランス精算の実施}：計画値と実績値の差分を精算する。
  \item \textbf{周波数維持・調整力確保}：電力系統の安定運用を担う。
\end{itemize}

新電力は送配電網を直接運用しないため、
一般送配電事業者の運用ルールに従う必要がある。

\subsubsection{（5）需給管理の経営的意義}

需給管理は、新電力の経営において以下の三つのリスクを左右する。

\begin{itemize}
  \item \textbf{インバランスリスク}：需給不一致による高額精算。
  \item \textbf{調達リスク}：JEPX の価格変動に伴う調達コストの急増。
  \item \textbf{信用リスク}：需給管理の失敗による事業撤退・倒産リスク。
\end{itemize}

特にインバランスリスクは、新電力の経営破綻の主要因として指摘されており、
需給管理の精度向上は事業継続の前提条件となっている。

以上のように、需給管理は新電力ビジネスの中核であり、
調達戦略と密接に連動しながら、事業の安定性を支える重要な機能である。

\subsection{販売：料金メニュー・付加価値サービス・顧客獲得戦略}

新電力の販売（小売）部門は、需要家に対して電力を供給し、収益を獲得する最終段階である。
販売戦略は、料金メニューの設計、付加価値サービスの提供、顧客獲得・維持のためのマーケティング活動など、
多様な要素から構成される。本節では、新電力の販売構造を体系的に整理する。

\subsubsection{（1）料金メニューの多様化：従量料金から環境価値メニューまで}

小売全面自由化により、新電力は旧一般電気事業者とは異なる料金メニューを自由に設計できるようになった。
料金メニューは、需要家の属性やニーズに応じて多様化しており、主に以下の種類が存在する。

\begin{itemize}
  \item \textbf{従量料金メニュー}：使用量に応じて料金が決まる最も一般的なメニュー。
  \item \textbf{時間帯別料金（TOU）}：昼夜やピーク・オフピークで料金を変動させるメニュー。
  \item \textbf{セット割メニュー}：ガス、通信、保険など他サービスとのセット販売。
  \item \textbf{再エネメニュー}：非化石証書を活用し、実質再エネ電力として提供するメニュー。
  \item \textbf{法人向け特別契約}：需要パターンに応じた個別最適化メニュー。
\end{itemize}

特に都市部では、環境価値を重視する需要家が増加しており、
再エネメニューの人気が高まっている。

\subsubsection{（2）付加価値サービス：差別化のための重要要素}

新電力は、料金だけでなく、付加価値サービスによって差別化を図る傾向が強い。
これは、家庭用市場では料金差が小さく、価格競争だけでは顧客獲得が難しいためである。

代表的な付加価値サービスは以下の通りである。

\begin{itemize}
  \item \textbf{ポイント還元}：T ポイント、Ponta、楽天ポイントなどとの連携。
  \item \textbf{スマートホーム連携}：IoT 機器と連動した省エネサービス。
  \item \textbf{環境価値サービス}：CO\textsubscript{2} 排出量の可視化、カーボンオフセット。
  \item \textbf{法人向けエネルギーマネジメント}：需要予測支援、ピークカット提案。
\end{itemize}

これらのサービスは、料金以外の価値を提供することで、
需要家のスイッチング行動を促す役割を果たす。

\subsubsection{（3）顧客獲得戦略：スイッチング行動と価格感応度への対応}

新電力の顧客獲得戦略は、需要家の行動特性を踏まえて設計される。
家庭用需要家は価格感応度が比較的低く、心理的コストが高いため、
料金差だけではスイッチングが進みにくい。

そのため、新電力は以下の戦略を組み合わせて顧客獲得を図る。

\begin{itemize}
  \item \textbf{料金比較サイトの活用}：価格差を可視化し、スイッチングの心理的障壁を下げる。
  \item \textbf{セット割による総合的メリットの提示}：通信・ガスとのセットで家計全体の節約を訴求。
  \item \textbf{環境価値の訴求}：再エネメニューを通じた環境配慮型ライフスタイルの提案。
  \item \textbf{地域密着型サービス}：地方では地域企業との連携が効果的。
\end{itemize}

法人需要家では、電気料金が直接的にコストに影響するため、
価格競争力が最も重要な要素となる。

\subsubsection{（4）販売部門の収益構造とリスク}

新電力の収益は、販売単価と調達コストの差（スプレッド）によって決まる。
しかし、販売部門には以下のリスクが存在する。

\begin{itemize}
  \item \textbf{調達コストの急騰}：JEPX 価格高騰時には販売単価を上回る調達コストが発生。
  \item \textbf{顧客離脱リスク}：料金改定やサービス低下により顧客が離脱する可能性。
  \item \textbf{信用リスク}：新電力の経営不安が需要家の信頼性評価に影響。
\end{itemize}

特に 2021 年の市場価格高騰時には、多くの新電力が販売単価を維持できず、
調達コストの急増により経営危機に陥った。

\subsubsection{（5）販売戦略の総括}

販売部門は、新電力の収益を直接生み出す最終段階であり、
料金メニュー、付加価値サービス、顧客獲得戦略の三つを統合的に設計する必要がある。
また、販売戦略は調達戦略・需給管理と密接に連動しており、
市場価格や制度変更に応じて柔軟に見直すことが求められる。

\subsection{新電力ビジネスの利害関係者：競争・協調・制度的制約}

新電力のビジネスモデルは、調達・需給管理・販売という三層構造の中で、
多様な利害関係者（ステークホルダー）が相互に関与する複雑な制度的枠組みの上に成り立っている。
本節では、図 1 のバリューチェーン全体を横断しながら、
新電力ビジネスに関与する主要主体とその利害関係を整理する。

\subsubsection{（1）発電事業者：競争相手であり供給源でもある二重の関係}

発電事業者は、新電力にとって電力調達の供給源であると同時に、
小売市場における競争相手でもある。
特に旧一般電気事業者は、発電部門で高いシェアを維持しており、
新電力はその供給力に依存せざるを得ない状況にある。

この関係は、以下のような利害構造を生む。

\begin{itemize}
  \item \textbf{協調関係}：相対契約を通じて安定的な電源供給を受ける。
  \item \textbf{競争関係}：小売市場では顧客獲得をめぐり直接競争する。
  \item \textbf{依存関係}：発電市場の寡占構造により、新電力は電源調達で不利な立場に置かれやすい。
\end{itemize}

このように、発電事業者は新電力にとって最も重要かつ複雑な利害関係を持つ主体である。

\subsubsection{（2）一般送配電事業者：中立的ネットワーク運用者としての役割}

一般送配電事業者は、送配電網の運用と需給調整を担う中立的主体であり、
新電力はそのネットワークを利用して需要家に電力を届ける。
送配電部門は法的分離により中立性が制度的に確保されているが、
新電力にとって以下のような制度的制約が存在する。

\begin{itemize}
  \item \textbf{託送料金制度}：新電力は送配電網の利用に対して託送料金を支払う必要がある。
  \item \textbf{接続ルール}：再エネ電源の接続制約や出力抑制が新電力の調達戦略に影響する。
  \item \textbf{需給実績の計測}：スマートメーターによる実績値計測はインバランス精算の基礎となる。
\end{itemize}

一般送配電事業者は競争主体ではないが、
その制度設計は新電力の事業運営に大きな影響を与える。

\subsubsection{（3）広域機関（OCCTO）：需給の番人としての制度的権限}

広域機関（OCCTO）は、需給計画の集約、広域的な需給調整、供給力評価、
インバランス精算ルールの策定など、電力システム全体の安定運用を担う。

新電力にとって、OCCTO は以下のような利害関係を持つ。

\begin{itemize}
  \item \textbf{制度的制約}：需給管理は OCCTO のルールに従う必要がある。
  \item \textbf{リスク管理}：インバランス精算は OCCTO の制度設計に依存する。
  \item \textbf{広域調整の恩恵}：需給逼迫時には広域的な調整により安定供給が確保される。
\end{itemize}

OCCTO は新電力の競争相手ではないが、
制度的枠組みを通じて新電力の経営リスクを左右する重要な主体である。

\subsubsection{（4）卸電力市場（JEPX）：価格形成の中心であり最大のリスク源}

JEPX は、新電力が電力を調達する主要市場であり、
スポット市場・先渡市場・時間前市場など複数の市場で構成される。

新電力にとっての利害関係は以下の通りである。

\begin{itemize}
  \item \textbf{機会}：市場価格が低い場合、調達コストを大幅に削減できる。
  \item \textbf{リスク}：市場価格が高騰した場合、調達コストが急増し経営危機に直結する。
  \item \textbf{依存}：相対契約だけでは不足する電力量を市場で補う必要がある。
\end{itemize}

特に 2021 年 1 月の価格高騰は、新電力の市場依存リスクを顕在化させた。

\subsubsection{（5）需要家：価格・サービス・信頼性を評価する最終主体}

需要家は、新電力の販売戦略の対象であり、
その行動は新電力の収益構造に直接影響する。

需要家の利害関係は以下の通りである。

\begin{itemize}
  \item \textbf{価格}：料金メニューの安さは重要だが、価格感応度は必ずしも高くない。
  \item \textbf{サービス}：ポイント還元、再エネメニュー、セット割など付加価値を重視する層が増加。
  \item \textbf{信頼性}：停電リスクや新電力の経営安定性に対する懸念がスイッチング行動を左右する。
\end{itemize}

需要家は新電力の競争環境を形成する主体であり、
その行動特性は販売戦略の前提条件となる。

\subsubsection{（6）行政（経済産業省）：制度設計者としての影響力}

経済産業省は、電力システム改革の制度設計を担う主体であり、
新電力の事業環境を規定する政策的権限を持つ。

行政の利害関係は以下の通りである。

\begin{itemize}
  \item \textbf{市場競争の促進}：自由化の目的は競争導入と効率化。
  \item \textbf{安定供給の確保}：市場原理だけでは安定供給が確保できない場合、制度的介入が必要。
  \item \textbf{脱炭素政策との整合性}：再エネ導入や非化石価値市場の設計に関与。
\end{itemize}

行政は競争政策と安定供給政策の両立を求められ、
その制度設計は新電力の事業戦略に直接影響する。

\subsubsection{（7）利害関係者構造の総括}

新電力ビジネスは、発電事業者、一般送配電事業者、OCCTO、JEPX、需要家、行政といった
多様な主体が相互に関与する複雑な制度的枠組みの中で成立している。
これらの主体は、競争・協調・制度的制約という三つの軸で関係を形成しており、
新電力はその中で調達戦略、需給管理、販売戦略を統合的に設計する必要がある。

この利害関係者構造は、新電力の経営リスクと事業機会を規定する基盤であり、
次章以降で検討する市場構造分析や倒産リスク分析の前提となる。

% ---------------------------------------------
% 1.4 卸電力市場の詳細
% ---------------------------------------------
\section{卸電力市場の詳細}

本節では、日本の卸電力市場（JEPX）を中心に、市場規模、価格推移、参加主体、
制度的要因などを定量データに基づいて整理する。
1.3 で示した新電力のビジネスモデルは、卸電力市場の価格変動や需給状況に強く依存しており、
市場の実態を把握することは新電力のリスク構造を理解する上で不可欠である。

\subsection{市場規模と取引量の推移（2020--2024）}

JEPX の年間取引量（TWh）、スポット市場のシェア、先渡市場の取引量などを示し、
自由化以降の市場拡大の実態を明らかにする。
また、取引量の季節性や再エネ出力との関連についても整理する。

\begin{figure}[H]
\centering
\includegraphics[width=14cm]{../figure/jepx_volume_2020_2024.pdf}
\caption{JEPX 年間取引量の推移（2020--2024）\\
{\small 出典：JEPX 公表データ（年間取引量統計）}}
\label{fig:jepx_volume}
\end{figure}



\subsection{市場参加者数の推移（新電力・発電事業者）}

小売電気事業者（新電力）、発電事業者、卸売事業者の登録数の推移を示し、
市場参入の動向と競争環境の変化を整理する。

\begin{figure}[H]
\centering
\includegraphics[width=14cm]{../figure/jepx_participants_2020_2024.pdf}
\caption{JEPX 市場参加者数の推移（2020--2024）\\
{\small 出典：JEPX 公表データ（市場参加者数統計）}}
\label{fig:jepx_participants}
\end{figure}

\subsection{退出企業数の推移（2020--2024）}

図 \ref{fig:exit_count} は、小売電気事業者の退出企業数の推移を示したものである。
2021 年以降、卸電力市場価格の高騰やインバランス料金の上昇を背景に、
退出企業数が急増していることが分かる。
退出には、法的倒産、実質的倒産、事業撤退、事業譲渡など複数の形態が含まれるが、
これらの分類および定義については次節（1.5 節）で詳述する。

\begin{figure}[H]
\centering
\includegraphics[width=14cm]{../figure/exit_count_2020_2024.pdf}
\caption{小売電気事業者の退出企業数の推移（2020--2024）\\
{\small 出典：経済産業省・電力広域的運営推進機関（OCCTO）公表資料を基に作成}}
\label{fig:exit_count}
\end{figure}


\subsection{スポット価格の推移と変動性}

スポット市場の平均価格、中央値、価格分布（ヒストグラム）、
および 2021 年 1 月の高騰事例を含む時系列推移を示す。
燃料価格、需給逼迫、再エネ出力の変動との関連も分析する。

\begin{figure}[H]
\centering
\includegraphics[width=15cm]{../figure/jepx_spot_timeseries_2020_2024.pdf}
\caption{JEPX スポット市場価格の推移（2020--2024）\\
{\small 出典：JEPX 公表データ（スポット市場価格統計）}}
\label{fig:jepx_spot_timeseries}
\end{figure}

\begin{figure}[H]
\centering
\includegraphics[width=14cm]{../figure/jepx_spot_hist_2020_2024.pdf}
\caption{JEPX スポット市場価格の分布（2020--2024）\\
{\small 出典：JEPX 公表データ（スポット市場価格統計）}}
\label{fig:jepx_spot_hist}
\end{figure}


\subsection{インバランス料金の推移}

インバランス料金の制度変更（2021 年 4 月）前後の比較、
高騰時のピーク値、季節性などを示し、
新電力の経営リスクとの関連を整理する。

\begin{figure}[H]
\centering
\includegraphics[width=15cm]{../figure/imbalance_timeseries_2020_2024.pdf}
\caption{インバランス料金の推移（2020--2024）\\
{\small 出典：OCCTO・電力広域的運営推進機関 公表データ}}
\label{fig:imbalance_timeseries}
\end{figure}


\subsection{需給逼迫アラート発令回数の推移}

2021 年以降の需給逼迫アラート（旧：需給ひっ迫注意報／警報）の発令回数を整理し、
需給逼迫が市場価格に与える影響を定量的に示す。
また、アラート発令日とスポット価格の関係も分析する。

\begin{figure}[H]
\centering
\includegraphics[width=14cm]{../figure/alert_count_2020_2024.pdf}
\caption{需給逼迫アラート発令回数の推移（2020--2024）\\
{\small 出典：経済産業省・電力広域的運営推進機関（OCCTO）公表資料}}
\label{fig:alert_count}
\end{figure}

\begin{figure}[H]
\centering
\includegraphics[width=14cm]{../figure/alert_vs_price_scatter.pdf}
\caption{アラート発令日とスポット価格の関係\\
{\small 出典：JEPX・OCCTO 公表データを基に作成}}
\label{fig:alert_vs_price_scatter}
\end{figure}


\subsection{JEPX の市場構造（スポット・先渡・時間前）}

各市場の役割、取引量の割合、価格形成への寄与を整理する。
また、先渡市場の流動性の低さや時間前市場の重要性についても説明する。

\begin{figure}[H]
\centering
\begin{tikzpicture}[node distance=2.2cm, >=stealth, thick]

% Nodes
\node[draw, rounded corners, fill=blue!10, minimum width=4cm, minimum height=1cm] (spot) {スポット市場};
\node[draw, rounded corners, fill=green!10, minimum width=4cm, minimum height=1cm, left=of spot] (fwd) {先渡市場（低流動性）};
\node[draw, rounded corners, fill=orange!10, minimum width=4cm, minimum height=1cm, right=of spot] (intraday) {時間前市場（調整）};

\node[draw, rounded corners, fill=gray!10, minimum width=4cm, minimum height=1cm, above=1.8cm of spot] (gen) {発電事業者};
\node[draw, rounded corners, fill=gray!10, minimum width=4cm, minimum height=1cm, below=1.8cm of spot] (retail) {小売電気事業者（新電力）};

% Arrows from generator
\draw[->] (gen) -- node[right]{供給計画} (spot);
\draw[->] (gen) -- node[above]{ヘッジ契約} (fwd);
\draw[->] (gen) -- node[above]{調整} (intraday);

% Arrows to retailer
\draw[->] (spot) -- node[right]{実需に応じた調達} (retail);
\draw[->] (fwd) -- node[above]{価格変動リスクの低減} (retail);
\draw[->] (intraday) -- node[above]{需給調整} (retail);

% Notes
\node[align=left, text width=6cm, below=0.8cm of retail] (note1)
{\small ・先渡市場は流動性が低く、十分なヘッジが困難\\
 ・時間前市場は需給調整の要だが、価格変動が大きい\\
 ・スポット市場への依存が高いほど、価格ショックに脆弱};

\end{tikzpicture}

\caption{JEPX の市場構造（スポット・先渡・時間前）\\
{\small 出典：電力広域的運営推進機関（OCCTO）資料を基に作成}}
\label{fig:jepx_market_structure}
\end{figure}

\subsection{価格形成メカニズム}

スポット市場の需給曲線、発電コスト構造、燃料価格、再エネ出力、
FIT 電源の市場放出量など、価格形成に影響を与える要因を体系的に整理する。

\begin{figure}[H]
\centering
\begin{tikzpicture}[scale=1.1, >=stealth, thick]

% Axes
\draw[->] (0,0) -- (7,0) node[right]{需要量};
\draw[->] (0,0) -- (0,6) node[above]{価格};

% Demand curve
\draw[blue, thick] (1,5) -- (6,1) node[midway, above right]{需要曲線};

% Supply curve (normal)
\draw[red, thick] (1,1) -- (6,5) node[midway, above left]{供給曲線};

% Supply curve shifted (renewables)
\draw[red!50, dashed, thick] (0.5,1) -- (5.5,5)
    node[midway, above left]{再エネ増加による供給曲線の左シフト};

% Intersection (equilibrium)
\filldraw[black] (3.5,3) circle (2pt);
\node[right] at (3.5,3) {均衡価格（スポット価格）};

% Cost blocks (baseload, mid, peak)
\draw[fill=green!20, opacity=0.6] (0.2,0.2) rectangle (2.5,0.8);
\node at (1.3,0.5) {\small ベースロード};

\draw[fill=yellow!20, opacity=0.6] (2.5,0.2) rectangle (4.5,1.2);
\node at (3.5,0.8) {\small ミドル電源};

\draw[fill=red!20, opacity=0.6] (4.5,0.2) rectangle (6.5,2.0);
\node at (5.5,1.4) {\small ピーク電源};

% FIT supply block
\draw[fill=blue!20, opacity=0.4] (0.2,1.2) rectangle (1.5,1.7);
\node at (0.85,1.45) {\tiny FIT電源};

\end{tikzpicture}

\caption{需給曲線の概念図（価格形成メカニズム）\\
{\small 出典：電力広域的運営推進機関（OCCTO）資料を基に作成}}
\label{fig:demand_supply_curve}
\end{figure}



\subsection{市場参加主体の役割}

電力市場には、発電事業者、小売電気事業者、卸売事業者、アグリゲーターなど、多様な主体が参加しており、それぞれが異なる目的と制約のもとで行動している。発電事業者は、燃料価格や設備稼働状況に応じて供給計画を策定し、市場に電力を供出する主体である。一方、小売電気事業者は、需要家との契約に基づき必要量を調達し、価格変動リスクを管理しながら販売を行う。卸売事業者は、需給調整やポートフォリオ最適化を目的として市場間で電力を仲介し、流動性の確保に寄与する。また、アグリゲーターは分散型電源や需要側リソースを束ね、需給調整市場や容量市場に参加することで、系統安定化に重要な役割を果たす。

これらの主体の行動は、スポット市場価格や需給逼迫リスクに直接的な影響を与える。特に、新電力は資本力や調達手段の制約からスポット市場への依存度が高く、他主体の行動変化が経営リスクに直結しやすい構造を持つ。したがって、市場参加主体の役割と相互作用を理解することは、価格形成メカニズムや倒産リスクの分析に不可欠である。

\subsection{制度的要因と市場への影響}

電力市場の価格形成や調達コストには、FIT 電源の市場放出、容量市場、需給調整市場など、制度的要因が大きく影響している。FIT 電源は固定価格買取制度のもとで導入が進んだ再生可能エネルギーであり、その市場放出量は供給曲線を左シフトさせ、スポット価格の低下圧力として作用する。一方、容量市場は将来の供給力確保を目的とした制度であり、発電事業者に対する固定的な収入源となるが、小売電気事業者にとっては追加的なコスト負担となる。

さらに、需給調整市場は系統安定化のための調整力を確保する仕組みであり、調整力の不足や制度変更に伴うコスト上昇は、小売電気事業者の調達コストを押し上げる要因となる。特に新電力は、調整力確保のためのリソースが限られているため、制度的コストの増加が経営に与える影響が相対的に大きい。

このように、制度的要因は単に市場価格に影響を与えるだけでなく、事業者間の競争構造やリスク分担のあり方を規定する重要な要素である。制度設計の変化が新電力の経営環境にどのような影響を及ぼすかを把握することは、倒産リスクの構造的理解に不可欠である。
。

以上のように、本節では日本の卸電力市場の制度的枠組みと価格形成の特徴を整理した。これらの市場構造は、新電力会社が直面する外部環境を規定するものであり、次節で扱う「新電力のリスク構造」を理解するための基盤となる。

% ---------------------------------------------
% 1.5 倒産・退出・継続の定義
% ---------------------------------------------

\section{倒産・退出・継続の定義}
\label{sec:definitions}

本研究では、新電力事業者の経営破綻や市場退出を分析するにあたり、
倒産研究の主要文献（Beaver, 1966; Altman, 1968; Ohlson, 1980; Shumway, 2001）で
用いられてきた概念的区分を尊重しつつ、
日本の倒産法制および電力市場の制度的特徴を踏まえた定義を採用する。
これにより、本研究の理論分析および実証分析において、
用語の意味が一貫して解釈されることを目的とする。

\subsection{法的倒産（Legal Bankruptcy）}

日本の倒産法制における法的倒産は、
破産法、民事再生法、会社更生法に基づく
裁判所主導の法的手続きが開始された状態を指す。

破産は清算型手続であり、事業継続を前提としない。
民事再生は中小企業向けの再建型手続であり、
経営陣が事業を継続しつつ再生計画の認可を目指す（DIP型）。
会社更生は大企業向けの再建型手続であり、
更生管財人が選任され、経営陣は原則として排除される。

Altman (1968)、Shumway (2001) などの倒産研究では、
これら三類型はいずれも「法的倒産」として扱われる。
本研究でも、官報、東京商工リサーチ、帝国データバンク等で
手続開始決定が確認された時点を法的倒産と定義する。


\subsection{実質的倒産（Substantive Bankruptcy / Failure）}

実質的倒産とは、法的手続きには至らないものの、
企業が事業継続不能な状態に陥った場合を指す。
Beaver (1966) や Ohlson (1980) では、債務不履行、資金繰り破綻、
営業停止などを含む広義の「failure」が用いられている。
新電力事業者の場合、供給停止、小売電気事業の廃止届、
債務不履行などが該当し、本研究ではこれらを「実質的倒産」と定義する。

\subsection{市場退出（Market Exit）}

市場退出とは、経営破綻に限らず、
事業撤退、事業譲渡、M\&A による吸収など、
企業が市場から退出する一連の事象を指す。
産業組織論（Geroski, 1995）では、退出は倒産とは区別されるが、
倒産リスクの顕在化の一形態として重要である。
本研究では、法的倒産・実質的倒産に加え、
事業譲渡や撤退も「市場退出」として整理する。

\subsubsection{事業撤退（Business Withdrawal）}

事業撤退とは、企業が法的倒産に至ることなく、
自主的に市場から退出する行為を指す。
新電力事業者の場合、経済産業省への小売電気事業廃止届の提出、
事業譲渡、M\&A による吸収合併などが該当する。

これらは財務破綻を伴わない場合も多く、
法的倒産とは区別されるべき退出形態である。
本研究では、事業撤退を「市場退出（exit）」の一類型として位置づける。

\subsection{事業継続（Going Concern）}

事業継続とは、企業が供給義務を果たしつつ、
財務的に継続可能な状態を指す。
会計基準（IAS 1）における「継続企業の前提」に基づき、
本研究では財務指標および市場リスク指標を用いて
事業継続可能性を評価する。


\subsection{定義の適用範囲と限界}

以上の定義は、新電力事業者の倒産リスクを分析するために
本研究が採用する概念的枠組みであるが、
すべての企業形態に機械的に適用できるわけではない。
特に、事業部制企業、多角化企業（コングロマリット）、
持株会社体制、連結決算ベースで複数事業を展開する企業では、
倒産・退出・継続の判断が事業単位と企業単位で一致しない場合がある。

\subsubsection{事業単位と企業単位の区別（Business Unit vs. Firm Level）}

新電力事業者は、小売電気事業という特定の事業単位で参入・退出する一方、
倒産は企業単位で発生する。
多角化企業（コングロマリット）や事業部制企業では、
特定事業の撤退が企業全体の倒産を意味しない場合があるため、
本研究では「事業単位の退出」と「企業単位の倒産」を明確に区別する。

\subsubsection{供給停止（Supply Suspension）}

供給停止とは、新電力事業者が電力供給を継続できなくなり、
顧客への供給義務を履行できない状態を指す。
法的倒産には該当しない場合があるが、
実質的倒産（failure）の重要な形態として本研究では位置づける。

\subsubsection{小売電気事業の廃止届（Business Termination Filing）}

小売電気事業者が経済産業省に提出する廃止届は、
事業撤退を示す公式な行政データである。
ただし、廃止届は必ずしも倒産を意味せず、
事業譲渡や戦略的撤退を含むため、
本研究では「市場退出（exit）」の一形態として扱う。

\begin{figure}[H]
\centering
\begin{tikzpicture}[
  scale=0.78,
  transform shape,
  node distance=12mm,
  box/.style={rectangle, draw, rounded corners, align=center, minimum width=32mm, minimum height=10mm}
]

% Top level
\node[box] (universe) {市場退出（Exit）};

% Second level
\node[box, below left=of universe] (legal) {法的倒産\\(破産・民事再生・会社更生)};
\node[box, below=of universe] (failure) {実質的倒産\\(供給停止・債務不履行)};
\node[box, below right=of universe] (withdraw) {事業撤退\\(廃止届・事業譲渡)};

% Third level (legal breakdown)
\node[box, below left=of legal] (hasan) {破産};
\node[box, below=of legal] (saisei) {民事再生};
\node[box, below right=of legal] (kosei) {会社更生};

% Going concern
\node[box, right=35mm of universe] (going) {事業継続\\(Going Concern)};

% Arrows
\draw[->] (universe) -- (legal);
\draw[->] (universe) -- (failure);
\draw[->] (universe) -- (withdraw);

\draw[->] (legal) -- (hasan);
\draw[->] (legal) -- (saisei);
\draw[->] (legal) -- (kosei);

\end{tikzpicture}
\caption{倒産・退出・継続の概念構造}
\label{fig:bankruptcy_structure}
\end{figure}


次節では、新電力市場における倒産リスクの制度的背景について述べる。

% ---------------------------------------------
% 1.6 新電力のリスク構造
% 新電力が直面する主要リスクを体系化し、倒産分析の基礎枠組みを提示する。
% ---------------------------------------------
\section{新電力のリスク構造}

本節では、新電力会社が直面する主要なリスクを体系的に整理し、本研究における倒産リスク分析の基礎枠組みを提示する。新電力の事業環境は、卸電力市場の価格変動、制度変更、需給管理、資金繰り、顧客行動といった多様な要因が複雑に絡み合う構造を持つ。そのため、個別の事象を列挙するだけではリスク全体像を把握することが難しく、倒産リスクの形成メカニズムを理解するためには、これらの要因を体系的に分類し、相互関係を明確にする必要がある。

本研究では、新電力のリスクを「市場リスク」「制度リスク」「オペレーショナルリスク」「財務リスク」「顧客リスク」の五つに分類する。この分類は恣意的なものではなく、国際的に広く採用されている企業リスク管理フレームワークである COSO ERM（Enterprise Risk Management）における四分類──すなわち（1）戦略リスク、（2）オペレーショナルリスク、（3）財務リスク、（4）コンプライアンス／制度リスク──と整合的である。

COSO ERM は、米国の企業会計・監査団体が共同で設立した COSO（Committee of Sponsoring Organizations of the Treadway Commission）が 2004 年に公表した企業リスク管理の国際的フレームワークであり、企業が直面するリスクを体系的に把握・管理するための標準的枠組みとして世界的に広く採用されている。COSO ERM では、企業リスクは外部環境の変化に起因する「戦略リスク」、業務プロセスやシステムに起因する「オペレーショナルリスク」、資金繰りや資本構成に関わる「財務リスク」、および法規制や制度変更に起因する「コンプライアンス／制度リスク」の四領域に分類される。この四分類は産業横断的に適用可能な汎用性を持ち、企業リスクを体系的に整理するための基礎的枠組みとして位置づけられている。

新電力における市場リスクは外部環境の変動が収益性に影響を与える点で戦略リスクに対応し、制度リスクは政策変更や制度設計の影響を受ける点でコンプライアンス／制度リスクに相当する。需給管理や業務運営に起因するオペレーショナルリスクは COSO の同名カテゴリーと一致し、資金繰りや資本構成に関わる財務リスクは財務リスクとして位置づけられる。

さらに、顧客リスクは、顧客の解約行動や未収金の発生といった外部要因であると同時に、企業内部の運営体制やキャッシュフロー管理にも影響を及ぼす点で、戦略リスクとオペレーショナルリスクの双方にまたがる複合的領域に位置づけられる。新電力の事業特性を踏まえると、顧客リスクを独立したカテゴリーとして扱うことは合理的であり、むしろ倒産リスクの分析に不可欠である。

また、倒産リスク研究の標準的枠組み（Altman、Ohlson、Shumway など）においても、企業の倒産は（1）外部ショック、（2）内部運営の脆弱性、（3）財務的耐性の三要因によって説明されるとされる。本研究で採用する五つのリスク分類は、この三要因モデルを新電力の制度・市場構造に適用した際の最適な分解であり、外部環境（市場・制度・顧客）、内部運営（オペレーション・顧客）、財務的耐性（財務）のすべてを網羅している。

以上より、本研究で整理する五つのリスク分類は、国際的な企業リスク管理フレームワークおよび倒産リスク研究の標準モデルと整合しつつ、新電力特有の市場構造と制度環境を反映したものであり、新電力の倒産リスクを包括的に分析するための必要かつ十分な枠組みである。図\ref{fig:chapter1_structure} に示すように、これらのリスクは相互に関連しながら企業の財務的耐性に影響を与え、最終的には倒産リスクとして顕在化する。本節では、以下で各リスクの内容とそのメカニズムを順に検討する。

\begin{figure}[htbp]
  \centering
  \begin{tikzpicture}[
    node distance=18mm,
    risk/.style={rectangle, draw, rounded corners, align=center, minimum width=24mm, minimum height=9mm},
    arrow/.style={->, thick},
    darrow/.style={->, thick, dashed}
  ]

    % 上段：4つのリスク
    \node[risk] (market) {市場リスク};
    \node[risk, right=of market] (institution) {制度リスク};
    \node[risk, right=of institution] (operational) {オペレーショナル\\リスク};
    \node[risk, right=of operational] (customer) {顧客リスク};

    % 中央下段：財務リスク
    \node[risk, below=18mm of institution] (financial) {財務リスク};

    % 右端：倒産リスク
    \node[risk, right=32mm of financial] (default) {倒産リスク};

    % 各リスクから倒産リスクへの矢印（直接の影響）
    \draw[arrow] (market.east) to[bend left=10] (default.north west);
    \draw[arrow] (institution.south east) to[bend left=5] (default.north);
    \draw[arrow] (operational.south) to[bend left=5] (default.north east);
    \draw[arrow] (customer.east) to[bend left=15] (default.east);
    \draw[arrow] (financial) -- (default);

    % 他リスクから財務リスクへの破線矢印（財務リスクの増幅・媒介）
    \draw[darrow] (market.south) -- (financial.north west);
    \draw[darrow] (institution.south) -- (financial.north);
    \draw[darrow] (operational.south west) -- (financial.north east);
    \draw[darrow] (customer.south west) to[bend right=15] (financial.east);

    % --- 凡例（実線・破線の意味） ---
    \node[align=left, anchor=north east] 
      at ($(current bounding box.south east)+(0,-0.5)$) {
        \begin{tabular}{ll}
        \raisebox{2pt}{\tikz{\draw[->, thick] (0,0) -- (0.8,0);}} 
          & 直接的な影響（倒産リスクに直接作用） \\
        \raisebox{2pt}{\tikz{\draw[->, thick, dashed] (0,0) -- (0.8,0);}} 
          & 財務リスクを通じた間接的な影響（増幅効果） \\
        \end{tabular}
      };

  \end{tikzpicture}
  \caption{新電力における5つのリスクと倒産リスクの関係}
  \label{fig:risk_network}
\end{figure}


\subsection{市場リスク：価格変動と需給逼迫}

市場リスクは、新電力が最も強く影響を受けるリスクである。スポット市場価格は、燃料価格、需給逼迫、再エネ出力の変動などにより大きく変動する。特に、2021年1月のような極端な価格高騰は、調達コストを急激に押し上げ、資本力の弱い新電力に深刻な影響を与える。市場リスクは短期的な収益変動に直結するだけでなく、各社の資本構成や財務的耐性がどの程度この変動に耐えられるかを左右する点で、倒産リスクとの関連が強い。

\subsection{制度リスク：インバランス制度・容量市場・FIT制度}

制度リスクは、政策変更や制度設計の影響を受けるリスクである。インバランス料金制度は需給管理の誤差に対して高額のペナルティを課す可能性があり、容量市場は将来の供給力確保を目的としつつも、新電力のコスト構造に影響を与える。また、FIT 電源の市場放出メカニズムは供給曲線の形状を変化させ、価格形成に影響する。これらの制度的要因は、単に市場価格に影響を与えるだけでなく、企業の資本負担や収益構造を通じて財務リスクを増幅させる点で、倒産リスクの形成に密接に関わる。

\subsection{オペレーショナルリスク：需給管理と業務運営}

需給管理の誤差はインバランス料金として精算されるため、新電力のオペレーショナルリスクの中心を成す。需要予測の精度、システム運用能力、アグリゲーション機能の有無などがリスクの大きさを左右する。また、業務運営の不備やシステム障害は、顧客対応や請求業務に影響を与え、キャッシュフローの悪化を招く可能性がある。これらのオペレーショナルリスクは、財務的余力の乏しい企業ほど影響が大きく、資本構成の違いが倒産リスクの差として顕在化しやすい。

\subsection{財務リスク：資金繰りと親会社支援}

新電力は、調達コストの変動やインバランス料金の急騰により、短期間で資金繰りが悪化する可能性がある。特に、資本力の弱い独立系新電力は、外部ショックに対する耐性が低い。一方、親会社の資本力や信用力が強い場合、内部資金移転（Internal Capital Markets）によってリスクが緩和される可能性がある\parencite{stein1997,almeida2004}。このように、財務リスクは市場リスクや制度リスクと相互作用しながら倒産リスクを形成するため、資本構成の違いが企業の生存可能性に大きな影響を与える。

\subsection{顧客リスク：解約・未収・需要変動}

顧客の解約行動は、新電力の収益構造に直接影響する。価格高騰時には、顧客が大手電力会社に戻る「逆スイッチング」が発生しやすい。また、未収金の増加はキャッシュフローを悪化させ、倒産リスクを高める。需要変動も、需給管理の誤差を通じてリスクを増幅させる。顧客リスクは市場要因と財務要因の双方に影響を及ぼすため、倒産リスクの包括的分析において重要な位置を占める。

% --- 章末の橋渡しパラグラフ（第2章への導入） ---
本節で整理した市場リスク、制度リスク、オペレーショナルリスク、財務リスク、顧客リスクは、新電力会社が直面する外部環境と内部構造の双方を反映したものである。これらのリスクは相互に関連しながら企業の資本構成や財務的耐性に影響を与え、最終的には倒産リスクとして顕在化する。本研究では、第2章以降において、これらの市場要因と財務要因を統合的に捉え、新電力会社の倒産リスクを包括的に分析する枠組みを構築する。



% ---------------------------------------------
% 1.6 本章のまとめ
% 第1章全体の論点を整理し、次章の倒産分析への接続を明確化する。
% ---------------------------------------------
\section{本章のまとめ}

本章では、世界のエネルギー情勢、日本の電力システム改革、新電力のビジネスモデル、卸電力市場の構造、そして新電力のリスク構造を体系的に整理した。まず、世界的なエネルギー転換の潮流と地政学リスクの高まりが、電力市場の価格変動性を増大させ、日本の電力市場にも直接的な影響を与えていることを確認した。次に、日本の電力システム改革の歴史的展開を整理し、2016年の小売全面自由化以降、新電力が多様な制度的枠組みの中で競争を行っている現状を明らかにした。

さらに、新電力のビジネスモデルを調達・需給管理・販売・回収の4機能に分解し、アセットライト型の構造が市場リスクや制度リスクを直接的に受けやすい特性を持つことを示した。卸電力市場（JEPX）の価格形成メカニズムや市場構造を整理し、特にスポット市場の変動性が新電力の経営に大きな影響を与えることを確認した。最後に、新電力のリスク構造を市場リスク、制度リスク、オペレーショナルリスク、財務リスク、顧客リスクの5類型に整理し、これらが相互に作用しながら事業者の財務健全性に影響を与えることを示した。

以上の整理を通じて、新電力の倒産リスクは単一の要因ではなく、外部環境、制度設計、内部能力、顧客行動といった複数の要素が複雑に絡み合う構造的問題であることが明らかとなった。次章では、本章で整理した制度的・市場的背景を踏まえ、新電力の倒産事例の分析および倒産リスクの定量的評価に進む。
